% Clase 14 --- El circuito RC y sus dos posiciones (§1).
% Movida acá desde el cuerpo de notes.tex al agregar el resto de las figuras,
% para que toda la clase use el mismo formato que el resto del curso.
% El interruptor es lo que hace las dos historias: en A la batería carga el
% condensador, en B queda descolgada y el condensador se descarga sobre R.
\begin{center}
\begin{circuitikz}[line width=0.9pt, color=figblue]
  \draw (0,0) to[battery1, l=$\varepsilon$] (0,3);
  \draw (0,3) to[switch, l=$A$, i=$I$] (3,3);
  \draw (3,3) to[R, l=$R$] (3,1.5);
  \draw (3,1.5) to[capacitor, l=$C$] (3,0);
  \draw (3,0) -- (0,0);
\end{circuitikz}\\[0.5em]
{\small\itshape Con el interruptor en \textbf{A} la batería queda conectada y el
condensador se carga; en \textbf{B} queda descolgada y lo que hay es un
condensador cargado descargándose sobre la resistencia. Conviene fijar la
convención de signos antes de escribir nada: se llama $Q$ a la carga de
\emph{una} placa e $I$ a la corriente entrante por ese lado, y con esa elección
$I = dQ/dt$. Si se hubiera llamado $Q$ a la carga de la otra placa quedaría
$I = -dQ/dt$, y todos los signos posteriores saldrían cambiados.}
\end{center}
